\chapter{Experiments}
\label{chap:experiments}

This chapter evaluates the empirical behavior of OGAD on a controlled collection of synthetic graph distributions and compares it with a direct graph diffusion model and classical probabilistic graph generators. The methodological components of the pipeline were introduced in Chapter~\ref{chap:method}; the purpose here is therefore to describe their concrete experimental instantiation, report the quantitative results, and analyze how consistently the compared models reproduce the structural properties of the reference graph distributions.

The experiments address three main questions. First, they investigate whether generation through an intermediate graphon representation can recover graph distributions characterized by different structural mechanisms. Second, they compare OGAD with DiGress and classical graph generators under a common evaluation procedure. Third, they examine whether relative performance is consistent across structural descriptors and datasets.

Five synthetic graph families are considered: Erd\H{o}s--R\'enyi (ER), Stochastic Block Model (SBM), Barab\'asi--Albert (BA), Watts--Strogatz (WS), and an expected-degree graph family (EDG). Generated graph distributions are evaluated through degree distribution, clustering-coefficient distribution, and normalized-Laplacian spectrum. For each descriptor, similarity to the held-out reference distribution is measured through Maximum Mean Discrepancy (MMD), following Section~\ref{sec:method_ggme}. Lower values indicate closer agreement with the reference graph distribution.

In addition to the raw MMD values, two rank-based aggregations are considered. The first averages model ranks across datasets independently for each structural descriptor, while the second averages ranks across the three descriptors independently for each dataset. These summaries provide a scale-independent view of relative performance.

Throughout this chapter, the term \emph{model} refers to one of the final methods included in the comparative evaluation, such as OGAD, DiGress, or a classical baseline. Within the OGAD hyperparameter search, the terminology follows Section~\ref{sec:method_d3pm_selection}: a \emph{D3PM candidate} denotes one candidate configuration, a \emph{retained checkpoint} is the best validation-loss checkpoint within that candidate, and GGME is subsequently used for generator selection across candidate configurations.


\section{Experimental Setup}
\label{sec:exp_setup}

\subsection{Datasets}
\label{subsec:exp_datasets}

The evaluation uses five synthetic datasets. Each contains \(300\) graphs with \(100\) nodes, partitioned into \(240\) training graphs and \(60\) test graphs. The final test set remains outside both training and candidate selection, following Section~\ref{sec:method_input_partitioning}, and is accessed only during final evaluation.

Table~\ref{tab:exp_datasets} summarizes the datasets.

\begin{table}[htbp]
    \centering
    \caption{Synthetic graph datasets used in the experimental evaluation.}
    \label{tab:exp_datasets}
    \begin{tabular}{llll}
        \hline
        Dataset & Configuration & Nodes & Main structural property \\
        \hline
        ER  & \texttt{er\_n100\_p005}     & 100 & homogeneous random connectivity \\
        SBM & \texttt{sbm\_n100\_4c}      & 100 & four-community block structure \\
        BA  & \texttt{ba\_n100\_m3}       & 100 & preferential attachment \\
        WS  & \texttt{ws\_n100\_k8\_p01} & 100 & small-world structure \\
        EDG & \texttt{expdeg\_n100\_lin}  & 100 & heterogeneous expected degrees \\
        \hline
    \end{tabular}
\end{table}

The suite covers substantially different forms of graph organization. ER represents homogeneous random connectivity and naturally corresponds to a constant graphon; SBM introduces block structure; and EDG introduces heterogeneous edge probabilities. BA and WS provide structurally different stress cases, characterized respectively by preferential attachment and by small-world organization with high clustering.

This diversity is particularly relevant to OGAD because its generative process operates through an intermediate graphon representation.


\subsection{Compared Models}
\label{subsec:exp_models}

OGAD is compared with DiGress and four classical probabilistic graph generators:

\begin{itemize}
    \item \textbf{DiGress}, a discrete diffusion model operating directly in the graph domain;
    \item \textbf{Erd\H{o}s--R\'enyi (ER)}, fitted through a global edge probability;
    \item \textbf{Configuration Model}, based on degree information estimated from the training graphs;
    \item \textbf{Stochastic Block Model (SBM)}, based on community-dependent connection probabilities;
    \item \textbf{Degree-Corrected Stochastic Block Model (DC-SBM)}, which additionally models degree heterogeneity within communities.
\end{itemize}

All models are evaluated through the same GGME procedure, so the reported differences arise from the generated graph distributions rather than from model-specific evaluation criteria.


\subsection{Baseline Configuration}
\label{subsec:exp_baseline_configuration}

The classical baselines are fitted exclusively from the training graphs and generate \(100\) graphs for final evaluation. The ER baseline estimates a global edge probability; the Configuration Model samples degree sequences from pooled training-set degree information; the SBM baseline estimates community structure and block-level connection probabilities; and DC-SBM additionally introduces heterogeneous node-level degree effects.

Classical baseline generation uses random seed \(42\).

DiGress is trained using the CommunityGraphs configuration. The executed configuration uses \(200\) epochs, batch size \(32\), gradient accumulation over \(16\) batches, \(1000\) diffusion steps, \(8\) graph-transformer layers, and \(8\) attention heads. The final DiGress graph set evaluated by GGME contains \(100\) graphs.

Where possible, DiGress and the classical baselines are generated once per dataset and reused across the corresponding OGAD runs. Thus, different OGAD configurations are compared against the same baseline graph collections.


\subsection{OGAD Experimental Configuration}
\label{subsec:exp_ogad_configuration}

The OGAD processing stages are defined in Chapter~\ref{chap:method}. Table~\ref{tab:exp_fixed_parameters} reports the principal numerical settings used in the successful experiments.

\begin{table}[htbp]
    \centering
    \caption{Main fixed parameters used in the OGAD experiments.}
    \label{tab:exp_fixed_parameters}
    \begin{tabular}{ll}
        \hline
        Parameter & Value \\
        \hline
        Graphon resolution & \(128\times128\) \\
        SIGL epochs & \(100\) \\
        SIGL learning rate & \(0.06\) \\
        SIGL INR batch size & \(1024\) \\
        Bootstrap replicates & \(500\) \\
        Bootstrap seed & \(0\) \\
        SIGL workers & \(4\) \\
        D3PM hidden dimension & \(512\) \\
        D3PM batch size & \(32\) \\
        D3PM diffusion steps & \(1000\) \\
        D3PM maximum epochs & \(300\) \\
        Early-stopping patience & \(25\) \\
        Early-stopping minimum improvement & \(5\times10^{-5}\) \\
        D3PM train/validation split & \(80/20\) \\
        D3PM split seed & \(0\) \\
        Generated graphon tensors per candidate & \(16\) \\
        Generated graphs per evaluation & \(100\) \\
        GGME seed & \(0\) \\
        \hline
    \end{tabular}
\end{table}

The main experimentally varied quantities are the number of clusters used before the SIGL bootstrap procedure and the hyperparameters defining each D3PM candidate configuration.

The principal sweep uses requested cluster counts

\[
C_{\mathrm{req}}\in\{1,4,8\}.
\]

As explained in Section~\ref{subsec:method_sigl_clustering}, the requested number of clusters can differ from the number of effective non-empty clusters obtained at runtime.

For each complete run in the main sweep, the D3PM candidate grid is defined by

\[
\eta
\in
\left\{
2\times10^{-5},
3\times10^{-5},
5\times10^{-5}
\right\},
\qquad
K\in\{10,14\},
\qquad
\lambda\in\{0.006,0.01\},
\]

resulting in twelve candidate configurations per OGAD run.

Additional follow-up experiments were performed for BA and WS using

\[
C_{\mathrm{req}}\in\{4,8,16\},
\]

and

\[
\eta
\in
\left\{
1\times10^{-5},
2\times10^{-5},
3\times10^{-5}
\right\},
\qquad
K\in\{10,14\},
\qquad
\lambda=0.006,
\]

resulting in six D3PM candidate configurations per run.


\subsection{Training and Candidate Selection}
\label{subsec:exp_candidate_selection}

The selection procedure follows the distinction introduced in Section~\ref{sec:method_d3pm_selection} between checkpoint selection within one D3PM candidate and generator selection across candidate configurations.

The \(240\) training graphs are first partitioned, using an \(80/20\) split with seed \(0\), into

\[
|\mathcal{G}_{\mathrm{model}}|=192
\]

graphs used by SIGL and

\[
|\mathcal{G}_{\mathrm{GGME-val}}|=48
\]

graphs reserved for graph-domain candidate validation.

The graphon collection produced by SIGL is independently divided into D3PM training and validation subsets, again using an \(80/20\) split with seed \(0\).

For each D3PM candidate configuration, training retains the checkpoint with the lowest validation loss. The retained checkpoint is then used to generate \(16\) graphon-like tensors, which are converted into finite graphs through the procedure described in Section~\ref{sec:method_graph_sampling}.

The graph distribution generated by each candidate is compared against \(\mathcal{G}_{\mathrm{GGME-val}}\) using GGME. The candidate obtaining the best GGME validation result according to the configured selection rule is retained as the selected D3PM candidate for that OGAD run.

The procedure can therefore be summarized as

\[
\begin{gathered}
\text{D3PM candidate}
\longrightarrow
\text{best checkpoint within candidate}
\longrightarrow
\text{generated graphs}
\\[0.4em]
\longrightarrow
\text{GGME validation}
\longrightarrow
\text{selected candidate}.
\end{gathered}
\]

The first selection operates in graphon representation space through the D3PM validation objective, whereas the second operates directly in graph space through GGME. The final test/reference graphs do not participate in either selection stage.


\subsection{Experimental Campaign and Completed Runs}
\label{subsec:exp_campaign}

The final experiments were obtained through a main broad sweep over all five datasets and a subsequent follow-up focused on BA and WS.

\begin{table}[htbp]
    \centering
    \caption{Main experimental sweeps contributing to the final evaluation.}
    \label{tab:exp_campaigns}
    \begin{tabular}{llll}
        \hline
        Sweep & Datasets & Cluster counts & D3PM learning rates \\
        \hline
        Main & ER, SBM, BA, WS, EDG
        & \(1,4,8\)
        & \(2,3,5\times10^{-5}\) \\
        Follow-up & BA, WS
        & \(4,8,16\)
        & \(1,2,3\times10^{-5}\) \\
        \hline
    \end{tabular}
\end{table}

Only runs completing all stages required for final GGME evaluation are included in the analysis.

\begin{table}[htbp]
    \centering
    \caption{Complete OGAD runs containing all three final GGME descriptor evaluations.}
    \label{tab:exp_complete_runs}
    \begin{tabular}{lc}
        \hline
        Dataset & Complete runs \\
        \hline
        ER  & 7 \\
        SBM & 6 \\
        BA  & 6 \\
        WS  & 4 \\
        EDG & 6 \\
        \hline
    \end{tabular}
\end{table}

Incomplete exploratory executions are excluded from the quantitative comparison.


\subsection{Evaluation Protocol}
\label{subsec:exp_evaluation}

Final evaluation follows Section~\ref{sec:method_ggme}. For every evaluated model, \(100\) generated graphs are compared with the held-out test graph distribution through:

\begin{enumerate}
    \item degree distribution (DD);
    \item clustering-coefficient distribution (CC);
    \item normalized-Laplacian spectrum (NL).
\end{enumerate}

For dataset \(D\), descriptor \(d\), and model \(m\), let

\[
M_{D,d}(m)
\]

denote the corresponding MMD value. Lower values indicate closer agreement between the generated and reference descriptor distributions.

A small number of reported MMD values are slightly negative. These values are retained exactly as emitted by the evaluation implementation, and ranking follows their reported numerical ordering.


\subsection{Aggregate Rank Analysis}
\label{subsec:exp_rank_method}

Raw MMD values cannot be meaningfully averaged across DD, CC, and NL because the descriptors operate on different representations and numerical scales. A rank-based analysis is therefore used as a complementary summary.

For every dataset \(D\), descriptor \(d\), and model \(m\), let

\[
r_{D,d}(m)\in\{1,\ldots,6\}
\]

be the rank obtained by sorting models by increasing MMD, with rank \(1\) assigned to the best result. Exact ties receive the average of the occupied ranks.

The first aggregation computes the mean rank across datasets for each descriptor:

\[
\overline{r}_{d}(m)
=
\frac{1}{5}
\sum_{D\in\mathcal{D}}
r_{D,d}(m),
\]

where

\[
\mathcal{D}
=
\{
\mathrm{ER},
\mathrm{SBM},
\mathrm{BA},
\mathrm{WS},
\mathrm{EDG}
\}.
\]

The second computes the mean rank across descriptors for each dataset:

\[
\overline{r}_{D}(m)
=
\frac{1}{3}
\sum_{d\in\{\mathrm{DD},\mathrm{CC},\mathrm{NL}\}}
r_{D,d}(m).
\]

Lower mean ranks are better in both cases.


\section{Quantitative Results}
\label{sec:exp_results}

\subsection{Erd\H{o}s--R\'enyi}
\label{subsec:exp_er}

\begin{table}[htbp]
    \centering
    \caption{MMD results on the Erd\H{o}s--R\'enyi dataset. Lower is better.}
    \label{tab:exp_er}
    \begin{tabular}{lrrr}
        \hline
        Model & DD & CC & NL \\
        \hline
        OGAD          & 0.080129 & 69.4112   & \(4.288\times10^{-6}\) \\
        DiGress       & 0.007058 & 21.5956   & \(\mathbf{-1.101\times10^{-6}}\) \\
        ER            & \textbf{0.001205} & \textbf{-4.56225} & \(5.754\times10^{-7}\) \\
        Configuration & 0.001266 & 7.37291   & \(-6.987\times10^{-7}\) \\
        SBM           & 0.004082 & 927.725   & \(2.563\times10^{-6}\) \\
        DC-SBM  & 0.161066 & 1291.50   & \(1.178\times10^{-5}\) \\
        \hline
    \end{tabular}
\end{table}

The matched ER baseline achieves the lowest degree and clustering discrepancies, while DiGress obtains the smallest reported spectral MMD. OGAD ranks \(5\), \(4\), and \(5\) for DD, CC, and NL, respectively. The result illustrates the advantage of a correctly specified probabilistic model when the target distribution directly matches its assumptions.


\subsection{Stochastic Block Model}
\label{subsec:exp_sbm}

\begin{table}[htbp]
    \centering
    \caption{MMD results on the Stochastic Block Model dataset. Lower is better.}
    \label{tab:exp_sbm}
    \begin{tabular}{lrrr}
        \hline
        Model & DD & CC & NL \\
        \hline
        OGAD          & 0.016365 & \textbf{0.203109} & \(4.035\times10^{-7}\) \\
        DiGress       & \textbf{0.002235} & 190.035 & \(1.005\times10^{-6}\) \\
        ER            & 0.003837 & 430.763 & \(\mathbf{-1.774\times10^{-7}}\) \\
        Configuration & 0.003710 & 614.534 & \(1.609\times10^{-8}\) \\
        SBM           & 0.014688 & 695.475 & \(1.962\times10^{-6}\) \\
        DC-SBM  & 0.217106 & 1394.32 & \(3.628\times10^{-6}\) \\
        \hline
    \end{tabular}
\end{table}

OGAD achieves its strongest result on clustering, with an MMD of \(0.203109\), substantially below the competing values and corresponding to rank \(1\). It ranks \(5\) on degree distribution and \(3\) on the normalized-Laplacian spectrum, showing that strong agreement on one structural property does not imply equally strong agreement on the others.


\subsection{Barab\'asi--Albert}
\label{subsec:exp_ba}

\begin{table}[htbp]
    \centering
    \caption{MMD results on the Barab\'asi--Albert dataset. Lower is better.}
    \label{tab:exp_ba}
    \begin{tabular}{lrrr}
        \hline
        Model & DD & CC & NL \\
        \hline
        OGAD          & 0.980363 & 508.541 & \(3.543\times10^{-5}\) \\
        DiGress       & 0.204514 & \textbf{78.0936} & \(\mathbf{1.979\times10^{-6}}\) \\
        ER            & 1.137100 & 1821.88 & \(1.337\times10^{-5}\) \\
        Configuration & \textbf{0.008423} & 900.171 & \(1.302\times10^{-5}\) \\
        SBM           & 1.135680 & 968.017 & \(1.524\times10^{-5}\) \\
        DC-SBM        & 0.882699 & 381.599 & \(1.777\times10^{-5}\) \\
        \hline
    \end{tabular}
\end{table}

BA is one of the most difficult settings for OGAD. The Configuration Model strongly outperforms it on degree distribution, while DiGress obtains the best clustering and spectral results. OGAD ranks \(4\), \(3\), and \(6\), yielding a mean rank of \(4.33\).


\subsection{Watts--Strogatz}
\label{subsec:exp_ws}

\begin{table}[htbp]
    \centering
    \caption{MMD results on the Watts--Strogatz dataset. Lower is better.}
    \label{tab:exp_ws}
    \begin{tabular}{lrrr}
        \hline
        Model & DD & CC & NL \\
        \hline
        OGAD          & 0.333326 & 2221.89 & \(1.290\times10^{-4}\) \\
        DiGress       & 0.290017 & \textbf{1505.24} & \(\mathbf{9.620\times10^{-5}}\) \\
        ER            & 0.322729 & 4532.93 & \(1.890\times10^{-4}\) \\
        Configuration & \textbf{0.040790} & 6003.34 & \(1.880\times10^{-4}\) \\
        SBM           & 0.297014 & 2783.35 & \(1.910\times10^{-4}\) \\
        DC-SBM        & 0.314288 & 2758.55 & \(1.860\times10^{-4}\) \\
        \hline
    \end{tabular}
\end{table}

OGAD ranks second for both clustering and normalized-Laplacian similarity, but sixth for degree distribution. The resulting mean rank is \(3.33\), showing that comparatively strong agreement on clustering and spectrum can coexist with a poor degree match.


\subsection{Expected-Degree Graphs}
\label{subsec:exp_edg}

\begin{table}[htbp]
    \centering
    \caption{MMD results on the expected-degree graph dataset. Lower is better.}
    \label{tab:exp_edg}
    \begin{tabular}{lrrr}
        \hline
        Model & DD & CC & NL \\
        \hline
        OGAD          & 0.022283 & \textbf{20.871} & \(1.696\times10^{-5}\) \\
        DiGress       & 0.020852 & 58.5847 & \(\mathbf{6.623\times10^{-6}}\) \\
        ER            & 0.519254 & 1397.30 & \(5.101\times10^{-5}\) \\
        Configuration & 0.022298 & 87.0578 & \(1.614\times10^{-5}\) \\
        SBM           & 0.355204 & 736.207 & \(1.326\times10^{-5}\) \\
        DC-SBM        & \textbf{0.012818} & 127.286 & \(5.929\times10^{-5}\) \\
        \hline
    \end{tabular}
\end{table}

EDG provides OGAD's strongest overall rank profile. OGAD ranks \(3\), \(1\), and \(4\) for DD, CC, and NL, respectively, giving a mean rank of \(2.67\). In particular, it obtains the best clustering result while remaining comparatively close to DiGress and the Configuration Model on degree distribution.


\section{Aggregate Rank Analysis}
\label{sec:exp_rank_analysis}

The raw MMD tables quantify the magnitude of descriptor-level discrepancies, while the rank analysis summarizes relative consistency across heterogeneous scales. The two views should therefore be interpreted jointly.


\subsection{Mean Rank Across Datasets}
\label{subsec:exp_rank_metric}

\begin{table}[htbp]
    \centering
    \caption{Mean rank across the five datasets, computed separately for each descriptor. Lower is better.}
    \label{tab:exp_rank_metric}
    \begin{tabular}{lccc}
        \hline
        Model & DD & CC & NL \\
        \hline
        OGAD          & 4.60 & 2.20 & 4.00 \\
        DiGress       & 2.20 & \textbf{1.80} & \textbf{1.60} \\
        ER            & 4.20 & 4.20 & 3.40 \\
        Configuration & \textbf{2.00} & 3.80 & 2.60 \\
        SBM           & 4.00 & 4.80 & 4.20 \\
        DC-SBM        & 4.00 & 4.20 & 5.20 \\
        \hline
    \end{tabular}
\end{table}

OGAD's strongest descriptor is clustering coefficient, with a mean rank of \(2.20\), second only to DiGress at \(1.80\). This summarizes its first-place results on SBM and EDG and its second-place result on WS.

Its mean degree rank is \(4.60\), compared with \(2.00\) for the Configuration Model and \(2.20\) for DiGress. For normalized-Laplacian spectrum, DiGress obtains the strongest average rank at \(1.60\), while OGAD reaches \(4.00\).

Clustering preservation therefore emerges as the most consistent empirical strength of OGAD.


\subsection{Mean Rank Across Structural Descriptors}
\label{subsec:exp_rank_dataset}

\begin{table}[htbp]
    \centering
    \caption{Mean rank across the three structural descriptors, computed separately for each dataset. Lower is better.}
    \label{tab:exp_rank_dataset}
    \begin{tabular}{lccccc}
        \hline
        Model & ER & SBM & BA & WS & EDG \\
        \hline
        OGAD          & 4.67 & 3.00 & 4.33 & 3.33 & 2.67 \\
        DiGress       & 2.67 & \textbf{2.33} & \textbf{1.33} & \textbf{1.33} & \textbf{1.67} \\
        ER            & \textbf{1.67} & \textbf{2.33} & 5.00 & 5.00 & 5.67 \\
        Configuration & 2.00 & 2.67 & 2.33 & 3.67 & 3.33 \\
        SBM           & 4.00 & 4.67 & 4.67 & 4.33 & 4.00 \\
        DC-SBM        & 6.00 & 6.00 & 3.33 & 3.33 & 3.67 \\
        \hline
    \end{tabular}
\end{table}

OGAD obtains its best dataset-level mean rank on EDG (\(2.67\)), followed by SBM (\(3.00\)) and WS (\(3.33\)). Its weakest aggregate results occur on BA and ER.

DiGress is the most consistent learned comparator, with mean ranks below \(3\) on every dataset. The matched ER baseline performs best on ER but deteriorates substantially on graph families generated through different mechanisms.


\section{Discussion}
\label{sec:exp_discussion}

The combined MMD and rank analyses show that generative performance depends strongly on both the target graph family and the structural property being evaluated.

\subsection{Clustering as the Strongest OGAD Descriptor}

The clearest empirical strength of OGAD concerns clustering-coefficient distributions. It achieves the lowest clustering MMD on SBM and EDG and ranks second on WS, producing a mean clustering rank of \(2.20\).

This pattern is particularly relevant to the graphon-mediated design of OGAD. SBM and EDG both admit natural descriptions through heterogeneous edge-probability structures, blockwise in the former case and smoother in the latter. Although the experiments do not isolate which pipeline stage causes this behavior, preservation of local and mesoscopic connectivity structure is the most consistent positive pattern observed for OGAD.


\subsection{Degree Distribution}

Degree distribution is instead OGAD's weakest descriptor, with a mean rank of \(4.60\).

The Configuration Model performs substantially better, with a mean rank of \(2.00\), and the contrast is especially pronounced on BA. This behavior is consistent with the fact that degree information is explicitly incorporated into the Configuration Model, whereas in OGAD node degrees emerge only after graphon estimation, diffusion generation, dequantization, latent-position sampling, and Bernoulli edge sampling.


\subsection{Normalized-Laplacian Spectrum}

DiGress is the most consistent model according to normalized-Laplacian similarity, with a mean rank of \(1.60\), compared with \(4.00\) for OGAD.

The contrast between OGAD's clustering and spectral results indicates that reproducing local or mesoscopic connectivity patterns does not necessarily imply equivalent preservation of global graph organization.


\subsection{Dependence on the Graph Family}

OGAD performs most favorably on EDG and SBM, with mean dataset ranks of \(2.67\) and \(3.00\), respectively. These are also the graph families whose structure admits a particularly direct interpretation through heterogeneous edge-probability functions.

BA is considerably more difficult, with a mean rank of \(4.33\), while ER favors its correctly specified classical baseline. WS occupies an intermediate position: OGAD is weak on degree distribution but ranks second on both clustering and spectrum.


\subsection{Comparison with DiGress}

DiGress is the most consistent learned comparator, with mean ranks of

\[
2.20,\qquad 1.80,\qquad 1.60
\]

for DD, CC, and NL, respectively.

The comparison is particularly informative because both methods rely on discrete diffusion but apply it to different representations. DiGress operates directly on graph structure, whereas OGAD applies diffusion to categorical graphon representations before sampling finite graphs.

The results suggest that these representations induce different structural biases: direct graph diffusion is more consistent overall, while graphon-mediated diffusion produces particularly strong clustering results on SBM and EDG.


\subsection{Role of Classical Baselines}

The classical baselines remain competitive when their structural assumptions align with either the target graph family or the evaluated statistic. The matched ER model obtains the strongest aggregate result on ER, while the Configuration Model achieves the best average degree rank.

Outside these favorable settings, however, their performance is less consistent. This highlights the complementary roles of specialized probabilistic generators and more general learned generative models.


\section{Experimental Limitations}
\label{sec:exp_limitations}

Several limitations should be considered when interpreting the results.

First, all experiments use synthetic graph families with \(100\) nodes, so the evaluation does not establish behavior on larger or real-world graph distributions.

Second, the search space is not perfectly homogeneous across datasets because BA and WS receive additional follow-up exploration, and the number of completed runs differs across graph families.

Third, DD, CC, and NL capture complementary but incomplete aspects of graph structure.

Fourth, mean rank removes the scale incompatibility between descriptors but discards information about the magnitude of MMD differences. The rank tables must therefore be interpreted together with the raw MMD values.

Finally, incomplete exploratory executions are excluded; the reported comparison includes only runs for which all three final GGME descriptor evaluations are available.


\section{Summary}
\label{sec:exp_summary}

The experiments show that OGAD exhibits a distinct descriptor- and dataset-dependent behavior rather than a uniform advantage over alternative generators.

Its strongest empirical result concerns clustering structure, for which it obtains a mean rank of \(2.20\) and the best MMD on both SBM and EDG. Its performance is weaker for degree distribution and normalized-Laplacian spectrum.

At dataset level, EDG and SBM are the most favorable settings for OGAD. DiGress is more consistent across the complete suite, while classical baselines remain particularly effective when their explicit assumptions match the target structure.

Overall, the results indicate that graphon-mediated generation introduces a distinct structural inductive bias: it is particularly effective for some clustering-related properties, while other aspects of graph structure remain better captured by direct graph diffusion or specialized probabilistic models.
